\begin{frame}
	\frametitle{Ответы на замечания Астапенко~В.А.}
	\textbf{Замечание}: Чем объясняется различие экспериментальной и теоретической кривой на рис.3.4?
	
	\textbf{Ответ}: Расхождение эксперименальной и модельной зависимости эффективности ГВГ от времени на рис.~3.4 наблюдается примерно до 39 часов. Возможная причина в том, что модель Антонюка предполагает постоянные параметры $u$ и $\alpha_2$, тогда как на начальном этапе существенны разброс времён релаксации ЭПЗ, неоднородное уширение и конкуренция переходов Ge-центр–матрица и Ge-центр–Ge-центр, что делает эти параметры зависящими от времени. В области насыщения (после 39 часов) влияние этих факторов ослабевает, и модель воспроизводит эксперимент, что видно по кинетике и КФС на рис. 3.5–3.6.
\end{frame}

\begin{frame}
	\frametitle{Ответы на замечания Астапенко~В.А.}
	\textbf{Замечание}: В Главе~1 приведен весьма полезный обзор литературы по теме диссертации, однако его относительный объем представляется излишне большим (46 стр. из 113 стр.)
	
	\textbf{Замечание}: Желательно было бы дать русскоязычное пояснению к термину «оптический полинг»
	
	\textbf{Ответ}: под термином «оптический полинг» понимается создание решётки квадратичной нелинейной восприимчивости при облучении ОВ излучением накачки.
\end{frame}

